\scanpage{133}
On prend syst. de copuiss. divisées sur $k[u]^+=S^+$ correspondant à celles \deleted{de}\inserted{canoniques} de $k'\{\{T\}\}^+$.

On prend $t\in k\{\{T\}\}$ donné par $t=T$, i.e.
\[
 t(u^n)=\left\{\begin{array}{ll}1&\text{si }n=1\\0&\text{si }n\neq1\end{array}\right.
\]
On a donc ici
\[
\begin{aligned}
 A_*(S,S^+,\gamma^{\unsure{*}},t)(k')
 &\simeq A_*(S^{k'},S^{+\,k'},\gamma_*^{k'},t^{k'})\\
 &\simeq A_*(k'\{\{T\}\},k'\{\{T\}\}^+,\gamma^{k'},T)\\
 &\simeq A_*(k')^\wedge
\end{aligned}
\]
On pose simplement
\[
 \deleted{\unsure{A_k^*}}\quad A_{k\unsure{*}}^*=A_k^*(S_k,\bar S_k,\gamma_k^*,t_k)
\]
de sorte que
\[
 A_*(k')^\wedge\simeq\operatorname{Hom}_k(A_k^*,k')
\]
``avec toutes les structures''. Notons qu'ici on est dans le cas où $S$, $\bar S$ et $\bar S\longrightarrow S$ sont gradués (graduation ordinaire de $k[u]$, correspondant à la graduation naturelle sur les $k'\{\{T\}\}$) et $t=T$ de \inserted{hom} degré 1, donc \deleted{$\unsure{A_k^*}$} $A_k^*$ admet une bigraduation. Celle-ci passe aux \deleted{\unsure{$C_*^*$}}
\[
 \deleted{C_*^{\mathrm{DR}}(X;k)}
\]
\[
 C_*^{\mathrm{DRpd}}(X,k)\overset{\mathrm{dfn}}\simeq C_*^{\mathrm{DR}}(X,k[u];k[u]^+,\gamma_k^*,t_k)
\]
\inserted{(avec les autres)}, et redonne par $\operatorname{Hom}_k(C_*^{\mathrm{DRpd}}(X,k),k')$ les structures connues (y compris bigraduation) sur le complexe de cochaînes de De Rham $C^*_{\mathrm{DRpd}}(X,k')$.

En fait tous les $C_*^{\mathrm{DRpd}}(X,k)$ ($k$ variable)
\margin{132}
